% GENERATED FILE. Edit canonical lessons, then run npm run generate:books.
% canonical-source-hash: fnv1a64:b2f5f92f47c06e85
% canonical-lessons: LA-C459-porticus, LA-C459-aedes, LA-C459-aedicula, LA-C459-pomerium, LA-C459-exstruo

\chapter{Colonnade, Temple, Shrine, Sacred boundary, To build up}
\label{ch:la-colonnade-temple-shrine-sacred-boundary-to-build-up}
\hlchaptermodality{459}

\begin{chapteropening}
\textbf{By the end of this chapter:} \emph{I can describe the public buildings of a town: a colonnade, a temple, a shrine, the sacred boundary, and putting buildings up.}

Complete the last lesson of chapter 459: I can describe the public buildings of a town: a colonnade, a temple, a shrine, the sacred boundary, and putting buildings up.
\end{chapteropening}

\section[porticus, porticūs]{porticus, porticūs — a colonnade, a covered walk, a portico}

\label{lesson:LA-C459-porticus}



\begin{quote}
Before the new one: say the Latin for a sewer, then the Latin for a water tower.
\end{quote}

\subsection*{You'll want to know: porticus, porticūs}
\textbf{porticus, porticūs} — \textquotedblleft{}a colonnade, a covered walk, a portico\textquotedblright{}.

\begin{quote}
Aestāte cīvēs in porticū ambulant, quod ibi umbra est.
\end{quote}

\begin{quote}
\emph{In summer the citizens walk in the colonnade, because there is shade there.}
\end{quote}

Fourth declension, feminine: \emph{porticus, porticūs f.} English \emph{portico} and \emph{porch} come from it.

\begin{cousinweb}
\emph{porticus} is built on \emph{porta} 'gate, passage' (EtymDB): a covered passage. English \emph{portico} comes from it through Italian, and \emph{porch} through French.
\end{cousinweb}

\begin{grammarlens}[title={What you've built}]
One more word for telling, explaining and describing at B1.
\end{grammarlens}

\subsection*{Your turn}
\begin{itemize}
\raggedright
  \item \emph{Say it:} \emph{porticus, porticūs}
  \item \emph{Say it:} \emph{porticus, porticūs}, once more
  \item \emph{Say it:} a sentence of your own with \emph{porticus, porticūs}
  \item \emph{Recall:} say the Latin for a sewer, then the Latin for a water tower, then say \emph{porticus, porticūs} again
\end{itemize}

\begin{tcolorbox}[breakable,colback=teal!4,colframe=teal!35!black,title={Before you move on}]
What does porticus, porticūs mean? (\textquotedblleft{}A colonnade, a covered walk, a portico\textquotedblright{}.) Say it once more.
\end{tcolorbox}
\section[aedēs, aedis]{aedēs, aedis — a temple; (plural) a house}

\label{lesson:LA-C459-aedes}



\begin{quote}
Before the new one: say the Latin for a water tower, then the Latin for a colonnade.
\end{quote}

\subsection*{You'll want to know: aedēs, aedis}
\textbf{aedēs, aedis} — \textquotedblleft{}a temple; (plural) a house\textquotedblright{}.

\begin{quote}
In forō aedēs magna et antīqua est.
\end{quote}

\begin{quote}
\emph{In the forum there is a great old temple.}
\end{quote}

Third declension, feminine: \emph{aedēs} or \emph{aedis, aedis f.} In the singular a temple, a house of a god; in the plural a house of people, with all its rooms.

\begin{cousinweb}
\emph{aedēs} is usually traced to PIE \emph{h\textsubscript{2}eyd\textsuperscript{h}-} 'to burn', since Lewis \& Short give 'hearth' as an early sense; \emph{aestās} 'summer' is from the same root. English \emph{edifice} comes from \emph{aedificium}, built on it.
\end{cousinweb}

\begin{grammarlens}[title={What you've built}]
One more word for telling, explaining and describing at B1.
\end{grammarlens}

\subsection*{Your turn}
\begin{itemize}
\raggedright
  \item \emph{Say it:} \emph{aedēs, aedis}
  \item \emph{Say it:} \emph{aedēs, aedis}, once more
  \item \emph{Say it:} a sentence of your own with \emph{aedēs, aedis}
  \item \emph{Recall:} say the Latin for a water tower, then the Latin for a colonnade, then say \emph{aedēs, aedis} again
\end{itemize}

\begin{tcolorbox}[breakable,colback=teal!4,colframe=teal!35!black,title={Before you move on}]
What does aedēs, aedis mean? (\textquotedblleft{}A temple; (plural) a house\textquotedblright{}.) Say it once more.
\end{tcolorbox}
\section[aedicula, aediculae]{aedicula, aediculae — a shrine, a small temple; a niche}

\label{lesson:LA-C459-aedicula}



\begin{quote}
Before the new one: say the Latin for a colonnade, then the Latin for a temple.
\end{quote}

\subsection*{You'll want to know: aedicula, aediculae}
\textbf{aedicula, aediculae} — \textquotedblleft{}a shrine, a small temple; a niche\textquotedblright{}.

\begin{quote}
In compitō aedicula parva est, et viātōrēs ibi flōrēs pōnunt.
\end{quote}

\begin{quote}
\emph{At the crossroads there is a small shrine, and travellers put flowers there.}
\end{quote}

First declension, feminine: \emph{aedicula, aediculae f.}, a diminutive of \emph{aedēs}. A little shrine, or a niche for a statue.

\begin{cousinweb}
\emph{aedicula} is a diminutive of \emph{aedēs} 'temple, house' (Lewis \& Short), made with \emph{-cula}. English keeps \emph{aedicule} in architecture.
\end{cousinweb}

\begin{grammarlens}[title={What you've built}]
One more word for telling, explaining and describing at B1.
\end{grammarlens}

\subsection*{Your turn}
\begin{itemize}
\raggedright
  \item \emph{Say it:} \emph{aedicula, aediculae}
  \item \emph{Say it:} \emph{aedicula, aediculae}, once more
  \item \emph{Say it:} a sentence of your own with \emph{aedicula, aediculae}
  \item \emph{Recall:} say the Latin for a colonnade, then the Latin for a temple, then say \emph{aedicula, aediculae} again
\end{itemize}

\begin{tcolorbox}[breakable,colback=teal!4,colframe=teal!35!black,title={Before you move on}]
What does aedicula, aediculae mean? (\textquotedblleft{}A shrine, a small temple; a niche\textquotedblright{}.) Say it once more.
\end{tcolorbox}
\section[pōmērium, pōmēriī]{pōmērium, pōmēriī — the sacred boundary (of a city)}

\label{lesson:LA-C459-pomerium}



\begin{quote}
Before the new one: say the Latin for a temple, then the Latin for a shrine.
\end{quote}

\subsection*{You'll want to know: pōmērium, pōmēriī}
\textbf{pōmērium, pōmēriī} — \textquotedblleft{}the sacred boundary (of a city)\textquotedblright{}.

\begin{quote}
Intrā pōmērium mīlitēs arma nōn gerunt.
\end{quote}

\begin{quote}
\emph{Inside the sacred boundary soldiers do not carry weapons.}
\end{quote}

Second declension, neuter: \emph{pōmērium, pōmēriī n.}, also \emph{pōmoerium}. The holy line round a Roman town, marked with stones.

\begin{cousinweb}
\emph{pōmērium} is usually explained, as by ancient writers, as \emph{post} 'behind' and \emph{moerus}, an old form of \emph{mūrus} 'wall' (Lewis \& Short): the strip of land along the wall. EtymDB lists it without a parent.
\end{cousinweb}

\begin{grammarlens}[title={What you've built}]
One more word for telling, explaining and describing at B1.
\end{grammarlens}

\subsection*{Your turn}
\begin{itemize}
\raggedright
  \item \emph{Say it:} \emph{pōmērium, pōmēriī}
  \item \emph{Say it:} \emph{pōmērium, pōmēriī}, once more
  \item \emph{Say it:} a sentence of your own with \emph{pōmērium, pōmēriī}
  \item \emph{Recall:} say the Latin for a temple, then the Latin for a shrine, then say \emph{pōmērium, pōmēriī} again
\end{itemize}

\begin{tcolorbox}[breakable,colback=teal!4,colframe=teal!35!black,title={Before you move on}]
What does pōmērium, pōmēriī mean? (\textquotedblleft{}The sacred boundary (of a city)\textquotedblright{}.) Say it once more.
\end{tcolorbox}
\section[exstruō, exstruere]{exstruō, exstruere — to build up, to raise, to pile up}

\label{lesson:LA-C459-exstruo}



\begin{quote}
Before the new one: say the Latin for a shrine, then the Latin for the sacred boundary.
\end{quote}

\subsection*{You'll want to know: exstruō, exstruere}
\textbf{exstruō, exstruere} — \textquotedblleft{}to build up, to raise, to pile up\textquotedblright{}.

\begin{quote}
Cīvēs novam porticum in forō exstruunt.
\end{quote}

\begin{quote}
\emph{The citizens put up a new colonnade in the forum.}
\end{quote}

\emph{exstruō, exstruere, exstrūxī, exstrūctum}. Also to pile up: \emph{mēnsās epulīs exstruere}, to load tables with food.

\begin{cousinweb}
\emph{exstruō} is built from \emph{ex-} 'up, out' and \emph{struō} 'to pile up, to build' (EtymDB; Lewis \& Short). English \emph{construct} and \emph{structure} come from \emph{struō}.
\end{cousinweb}

\begin{grammarlens}[title={What you've built}]
One more word for telling, explaining and describing at B1.
\end{grammarlens}

\subsection*{Your turn}
\begin{itemize}
\raggedright
  \item \emph{Say it:} \emph{exstruō, exstruere}
  \item \emph{Say it:} \emph{exstruō, exstruere}, once more
  \item \emph{Say it:} a sentence of your own with \emph{exstruō, exstruere}
  \item \emph{Recall:} say the Latin for a shrine, then the Latin for the sacred boundary, then say \emph{exstruō, exstruere} again
\end{itemize}

\begin{tcolorbox}[breakable,colback=teal!4,colframe=teal!35!black,title={Before you move on}]
What does exstruō, exstruere mean? (\textquotedblleft{}To build up, to raise, to pile up\textquotedblright{}.) Say it once more.
\end{tcolorbox}
